\section{Representation and Editability}
\label{sec:representation}
We treat the \emph{delivered artifact}, rather than the architecture of its generator, as the first classification axis. In design production the image rendered to viewers and the editable artifact opened by designers are distinct observables. Table~\ref{tab:representation} summarizes common affordances and the evidence that is still needed.

% Independent source table.
\begin{table*}[t]
\centering
\caption{Representation--editability crosswalk.
The listed operations are possible affordances, not guarantees of empirically validated editing.}
\label{tab:representation}
\small
\begin{tabularx}{\textwidth}{@{}p{0.18\textwidth}XX@{}}
\toprule
\textbf{Delivered artifact} & \textbf{Potential post-generation operation} & \textbf{Critical evidence boundary} \\
\midrule
Layout boxes and classes & Move, resize, or constrain element geometry. & Geometry does not establish editable text, asset identities, or a portable design file. \\
Flattened raster image & Pixel-level modification, inpainting, or mask-based regeneration. & Raster image editing is valid, but it does not imply native vector or text-object editing. \\
SVG vector paths & Modify paths, fill, stroke, and layer order; interpolate shapes. & Outlined letters are not semantic font objects; path access is not an editing-success score. \\
Layered image/document & Reposition, replace, or re-render separable layers and components. & RGBA layers may still contain fused objects; alpha channels alone do not ensure semantic separation. \\
Code-native HTML/CSS or presentation objects & Edit DOM/object nodes, text, and layout rules; re-render and inspect. & Syntactically valid code need not preserve intended layout or succeed at natural-language edit requests. \\
\bottomrule
\end{tabularx}
\end{table*}


\subsection{Layout Geometry and Raster Composites}
LayoutDM models category and quantized spatial attributes through discrete diffusion with conditional generation from partial layouts \citep{layoutdm2023}. LayoutFlow models continuous geometric generation \citep{layoutflow2024}. Such systems provide useful structural arrangements, but layout coordinates need not include the image and text assets needed for a fully editable poster. Content-aware layout systems additionally place foreground elements in visual context \citep{cglgan2022,layoutdetr2024}.

Raster generation delivers an integrated rendering. A clean appearance is valuable, but targeted changes and preservation of neighboring content constitute a separate image-editing problem. SAWNA exemplifies a boundary: it enforces user-reserved blank regions during image synthesis to support later composition \citep{sawna2025}. Those regions are useful but do not establish post-export semantic object manipulation.

\subsection{Vector Path and Layer Structures}
DeepSVG models drawing commands and studies reconstruction and latent interpolation \citep{deepsvg2020}; SVGformer learns representations of continuous vector graphics \citep{svgformer2023}. VectorFusion optimizes differentiable SVG paths with text-to-image diffusion priors \citep{vectorfusion2023}. IconShop models SVG icon paths and demonstrates editing by path-level infilling \citep{iconshop2023}. The existence of vector paths does not establish instruction-driven edit accuracy or unchanged-object preservation.

Hierarchical systems including COLE and OpenCOLE generate designs through intermediate composition steps \citep{cole2023,opencole2024}. LayerD reconstructs layered structure from flattened designs \citep{layerd2025}. Code-native systems represent text and assets with manipulable HTML/CSS elements \citep{designascode2026,editablevisual2026}. Their artifact format provides substantial leverage for editing, yet success under a controlled editing protocol remains an independent empirical question.

\subsection{Four Types of Evidence}
We distinguish: \emph{format affordance}, where structure permits a type of manipulation; \emph{workflow demonstration}, where authors show specific edits; \emph{human editing efficiency}, where users are measured completing a design task; and \emph{quantitative edit-operation evaluation}, where target change and unintended changes are scored. The levels are not a universal ordinal quality scale. A qualitative demonstration can be convincing for one operation while failing to establish statistical reliability, and a controlled human study can measure a different property from automatic instruction following.
